
\documentclass[lineno]{jfm}

% Build fix only -- no content. jfm.cls defers \usepackage{lineno-FLM} to
% \AtBeginDocument, but LaTeX reads paper.aux at \begin{document} before that
% hook fires, so the \@LN@col written into the aux by the previous pdflatex
% pass is undefined when it is read back and the build dies with "Undefined
% control sequence". Providing a no-op lets the aux be read; lineno redefines
% it properly a moment later.
\makeatletter
\providecommand\@LN@col[1]{}
\makeatother

\usepackage{graphicx}
%\usepackage{epstopdf,epsfig}
\usepackage{newtxtext}
\usepackage{newtxmath}
\usepackage{natbib}
\usepackage{hyperref}
\hypersetup{
    colorlinks = true,
    urlcolor   = blue,
    citecolor  = black,
}
% Provenance markers: every value injected from the project is colored, and
% every injected value, figure and question block is logged into
% main.provenance.json for the editor and the hub to follow back to its
% source. Drop the option for the journal version.
\usepackage[provenance]{calkit}
\newtheorem{lemma}{Lemma}
\newtheorem{corollary}{Corollary}
\newcommand{\RomanNumeralCaps}[1]
\linenumbers

% {\MakeUppercase{\romannumeral #1}}

\title{Data-driven closure discovery for bypass transition yields
non-transferable coefficients}

\author{P. Bachant\aff{1}
  \corresp{\email{pbachant@caltech.edu}}}

\affiliation{\aff{1}California Institute of Technology, Pasadena, CA 91125, USA}

\makeatletter%
\newcommand\result[1][all]{%
  \ifnum\pdfstrcmp{#1}{all}=0%
    \def\result@out{%
    \{%
      "CepsJimenez": 7.681, %
      "CepsPre": 1.61, %
      "CepsRatio": 4.3, %
      "CepsTurb": 6.973, %
      "CepsTurbSd": 0.278, %
      "ClipNewCf": 0.235, %
      "ClipNewU": 0.0389, %
      "CollapseGamma": 64.5, %
      "CollapseGammaCorr": 0.9832, %
      "CollapseTurnovers": 52.9, %
      "CollapseX": 60.1, %
      "DeltaInlet": 1.03, %
      "EpsDownToUp": {-}0.87, %
      "EpsInSample": 0.765, %
      "EpsUpToDown": 0.09, %
      "FitNoiseSd": 2.09, %
      "JimenezSelfFit": 0.962, %
      "LMNewCf": 0.173, %
      "LMNewU": 0.0337, %
      "LMOldCf": 0.701, %
      "LMOldU": 0.0501, %
      "LamCMax": 558.0, %
      "LamCMin": 445.0, %
      "LaminarNewCf": 0.51, %
      "NSeeds": 4, %
      "NTerms": 8, %
      "NVariants": 3, %
      "PoverEpsPeak": 2.58, %
      "PoverEpsPeakX": 110.0, %
      "ReThetatInlet": 210.0, %
      "SSTNewCf": 0.718, %
      "StressCoeffDisagree": 1.0, %
      "StressDiagnostic": 0.971, %
      "StressDownToUp": {-}3394.0, %
      "StressInSample": 0.953, %
      "StressToJimenez": {-}472.0, %
      "TuInlet": 2.54, %
      "XInlet": 30.2, %
      "XOutlet": 1000.1, %
      "YMax": 26.5, %
      "aOneJimenez": 0.1311, %
      "aOneJimenezSd": 0.0009, %
      "aOneRelDiff": 5.0, %
      "aOneTransitional": 0.137%
    \}}%
  \else%
    \ifnum\pdfstrcmp{#1}{CepsJimenez}=0%
      \def\result@out{%
    7.681}%
    \else%
      \ifnum\pdfstrcmp{#1}{CepsPre}=0%
      \def\result@out{%
    1.61}%
    \else%
      \ifnum\pdfstrcmp{#1}{CepsRatio}=0%
      \def\result@out{%
    4.3}%
    \else%
      \ifnum\pdfstrcmp{#1}{CepsTurb}=0%
      \def\result@out{%
    6.973}%
    \else%
      \ifnum\pdfstrcmp{#1}{CepsTurbSd}=0%
      \def\result@out{%
    0.278}%
    \else%
      \ifnum\pdfstrcmp{#1}{ClipNewCf}=0%
      \def\result@out{%
    0.235}%
    \else%
      \ifnum\pdfstrcmp{#1}{ClipNewU}=0%
      \def\result@out{%
    0.0389}%
    \else%
      \ifnum\pdfstrcmp{#1}{CollapseGamma}=0%
      \def\result@out{%
    64.5}%
    \else%
      \ifnum\pdfstrcmp{#1}{CollapseGammaCorr}=0%
      \def\result@out{%
    0.9832}%
    \else%
      \ifnum\pdfstrcmp{#1}{CollapseTurnovers}=0%
      \def\result@out{%
    52.9}%
    \else%
      \ifnum\pdfstrcmp{#1}{CollapseX}=0%
      \def\result@out{%
    60.1}%
    \else%
      \ifnum\pdfstrcmp{#1}{DeltaInlet}=0%
      \def\result@out{%
    1.03}%
    \else%
      \ifnum\pdfstrcmp{#1}{EpsDownToUp}=0%
      \def\result@out{%
    {-}0.87}%
    \else%
      \ifnum\pdfstrcmp{#1}{EpsInSample}=0%
      \def\result@out{%
    0.765}%
    \else%
      \ifnum\pdfstrcmp{#1}{EpsUpToDown}=0%
      \def\result@out{%
    0.09}%
    \else%
      \ifnum\pdfstrcmp{#1}{FitNoiseSd}=0%
      \def\result@out{%
    2.09}%
    \else%
      \ifnum\pdfstrcmp{#1}{JimenezSelfFit}=0%
      \def\result@out{%
    0.962}%
    \else%
      \ifnum\pdfstrcmp{#1}{LMNewCf}=0%
      \def\result@out{%
    0.173}%
    \else%
      \ifnum\pdfstrcmp{#1}{LMNewU}=0%
      \def\result@out{%
    0.0337}%
    \else%
      \ifnum\pdfstrcmp{#1}{LMOldCf}=0%
      \def\result@out{%
    0.701}%
    \else%
      \ifnum\pdfstrcmp{#1}{LMOldU}=0%
      \def\result@out{%
    0.0501}%
    \else%
      \ifnum\pdfstrcmp{#1}{LamCMax}=0%
      \def\result@out{%
    558.0}%
    \else%
      \ifnum\pdfstrcmp{#1}{LamCMin}=0%
      \def\result@out{%
    445.0}%
    \else%
      \ifnum\pdfstrcmp{#1}{LaminarNewCf}=0%
      \def\result@out{%
    0.51}%
    \else%
      \ifnum\pdfstrcmp{#1}{NSeeds}=0%
      \def\result@out{%
    4}%
    \else%
      \ifnum\pdfstrcmp{#1}{NTerms}=0%
      \def\result@out{%
    8}%
    \else%
      \ifnum\pdfstrcmp{#1}{NVariants}=0%
      \def\result@out{%
    3}%
    \else%
      \ifnum\pdfstrcmp{#1}{PoverEpsPeak}=0%
      \def\result@out{%
    2.58}%
    \else%
      \ifnum\pdfstrcmp{#1}{PoverEpsPeakX}=0%
      \def\result@out{%
    110.0}%
    \else%
      \ifnum\pdfstrcmp{#1}{ReThetatInlet}=0%
      \def\result@out{%
    210.0}%
    \else%
      \ifnum\pdfstrcmp{#1}{SSTNewCf}=0%
      \def\result@out{%
    0.718}%
    \else%
      \ifnum\pdfstrcmp{#1}{StressCoeffDisagree}=0%
      \def\result@out{%
    1.0}%
    \else%
      \ifnum\pdfstrcmp{#1}{StressDiagnostic}=0%
      \def\result@out{%
    0.971}%
    \else%
      \ifnum\pdfstrcmp{#1}{StressDownToUp}=0%
      \def\result@out{%
    {-}3394.0}%
    \else%
      \ifnum\pdfstrcmp{#1}{StressInSample}=0%
      \def\result@out{%
    0.953}%
    \else%
      \ifnum\pdfstrcmp{#1}{StressToJimenez}=0%
      \def\result@out{%
    {-}472.0}%
    \else%
      \ifnum\pdfstrcmp{#1}{TuInlet}=0%
      \def\result@out{%
    2.54}%
    \else%
      \ifnum\pdfstrcmp{#1}{XInlet}=0%
      \def\result@out{%
    30.2}%
    \else%
      \ifnum\pdfstrcmp{#1}{XOutlet}=0%
      \def\result@out{%
    1000.1}%
    \else%
      \ifnum\pdfstrcmp{#1}{YMax}=0%
      \def\result@out{%
    26.5}%
    \else%
      \ifnum\pdfstrcmp{#1}{aOneJimenez}=0%
      \def\result@out{%
    0.1311}%
    \else%
      \ifnum\pdfstrcmp{#1}{aOneJimenezSd}=0%
      \def\result@out{%
    0.0009}%
    \else%
      \ifnum\pdfstrcmp{#1}{aOneRelDiff}=0%
      \def\result@out{%
    5.0}%
    \else%
      \ifnum\pdfstrcmp{#1}{aOneTransitional}=0%
      \def\result@out{%
    0.137}%
    \else%
      \def\result@out{%
    ??}%
    \fi\fi\fi\fi\fi\fi\fi\fi\fi\fi\fi\fi\fi\fi\fi\fi\fi\fi\fi\fi\fi\fi\fi\fi\fi\fi\fi\fi\fi\fi\fi\fi\fi\fi\fi\fi\fi\fi\fi\fi\fi\fi\fi\fi
  \fi
  \result@out
}%
\makeatother%
%% Generated by Calkit. Do not edit.
%% Values are wrapped in \ckvalue so calkit.sty can mark and log where
%% they came from; without the package they print as plain text.
\providecommand\ckvalue[4]{#2}%
\providecommand\ckblock[2]{}%
\makeatletter%
\newcommand\ckquestion[1][all]{%
  \ifnum\pdfstrcmp{#1}{all}=0%
    \def\ckquestion@out{%
      1, 2, 3, 4, 5, 6, 7, 8, 9, 10, 11, 12, 13}%
  \else\ifnum\pdfstrcmp{#1}{1}=0%
    \def\ckquestion@out{%
      Can new terms in the steady RANS equations make the mean flow match the DNS?}%
  \else\ifnum\pdfstrcmp{#1}{2}=0%
    \def\ckquestion@out{%
      Can the Reynolds stress be modeled as a function of the mean flow?}%
  \else\ifnum\pdfstrcmp{#1}{3}=0%
    \def\ckquestion@out{%
      Is the clipping analogy - a flow holding more energy than smooth waves can carry, so it clips and redistributes into higher harmonics - load-bearing, or decorative?}%
  \else\ifnum\pdfstrcmp{#1}{4}=0%
    \def\ckquestion@out{%
      Which local quantity should carry the transition threshold?}%
  \else\ifnum\pdfstrcmp{#1}{5}=0%
    \def\ckquestion@out{%
      Does transition redistribute energy between velocity components, or create correlation between them?}%
  \else\ifnum\pdfstrcmp{#1}{6}=0%
    \def\ckquestion@out{%
      Can entropy alone serve as the closure variable?}%
  \else\ifnum\pdfstrcmp{#1}{7}=0%
    \def\ckquestion@out{%
      If entropy cannot be rejected as fast as the wall injects exergy, is the excess stored in the turbulence?}%
  \else\ifnum\pdfstrcmp{#1}{8}=0%
    \def\ckquestion@out{%
      Can a search discover the closure structure rather than have us impose it?}%
  \else\ifnum\pdfstrcmp{#1}{9}=0%
    \def\ckquestion@out{%
      Is the parabolic screening solver accurate enough to rank closures against each other?}%
  \else\ifnum\pdfstrcmp{#1}{10}=0%
    \def\ckquestion@out{%
      Does the closure survive a flow with no wall?}%
  \else\ifnum\pdfstrcmp{#1}{11}=0%
    \def\ckquestion@out{%
      Does the transition threshold depend on the free-stream turbulence level, as it physically must?}%
  \else\ifnum\pdfstrcmp{#1}{12}=0%
    \def\ckquestion@out{%
      Can we write an evolution equation for a structure variable such as Q or lambda\_2, rather than for an energy, from an advection-diffusion skeleton plus coupling terms from the mean momentum and continuity equations?}%
  \else\ifnum\pdfstrcmp{#1}{13}=0%
    \def\ckquestion@out{%
      Do the Reynolds stress transport terms proposed in notebooks/main.ipynb outperform the activation approach?}%
  \else\PackageError{calkit}{Unknown key '#1' for \string\ckquestion}{}%
  \fi\fi\fi\fi\fi\fi\fi\fi\fi\fi\fi\fi\fi\fi%
  \ckquestion@out}%
\makeatother%
\makeatletter%
\newcommand\ckhypothesis[1][all]{%
  \ifnum\pdfstrcmp{#1}{all}=0%
    \def\ckhypothesis@out{%
      1, 2, 3, 5, 6, 7, 10, 11}%
  \else\ifnum\pdfstrcmp{#1}{1}=0%
    \def\ckhypothesis@out{%
      A threshold, not extra smooth terms, is needed, because no k-epsilon coefficient set makes the laminar state a fixed point of the model.}%
  \else\ifnum\pdfstrcmp{#1}{2}=0%
    \def\ckhypothesis@out{%
      A scalar eddy viscosity relating the Reynolds stress to the mean strain is sufficient.}%
  \else\ifnum\pdfstrcmp{#1}{3}=0%
    \def\ckhypothesis@out{%
      Decorative; the model works because of its other terms.}%
  \else\ifnum\pdfstrcmp{#1}{5}=0%
    \def\ckhypothesis@out{%
      Transition redistributes energy from the streamwise into the wall-normal and spanwise components.}%
  \else\ifnum\pdfstrcmp{#1}{6}=0%
    \def\ckhypothesis@out{%
      Transition can be modeled as entropy production in the fluctuation field, so a transported entropy would close the system.}%
  \else\ifnum\pdfstrcmp{#1}{7}=0%
    \def\ckhypothesis@out{%
      Turbulence buffers exergy injection against entropy rejection, and the buffering matters most where the two are most out of balance.}%
  \else\ifnum\pdfstrcmp{#1}{10}=0%
    \def\ckhypothesis@out{%
      A closure whose length scale and transition gate are built from wall distance produces nothing in a free shear layer; the transported-omega variant, whose length scale is k/omega, does not need the wall.}%
  \else\ifnum\pdfstrcmp{#1}{11}=0%
    \def\ckhypothesis@out{%
      Higher free-stream turbulence lowers the critical Reynolds number.}%
  \else\PackageError{calkit}{Unknown key '#1' for \string\ckhypothesis}{}%
  \fi\fi\fi\fi\fi\fi\fi\fi\fi%
  \ckhypothesis@out}%
\makeatother%
\makeatletter%
\newcommand\ckanswer[1][all]{%
  \ifnum\pdfstrcmp{#1}{all}=0%
    \def\ckanswer@out{%
      1, 2, 3, 4, 5, 6, 7, 8, 9, 10}%
  \else\ifnum\pdfstrcmp{#1}{1}=0%
    \def\ckanswer@out{%
      Yes. A transported activation fraction with a rectified (clipping) source added to k-omega cuts the mean-velocity error against the DNS by \ckvalue{improvement}{5.1}{results/findings.json}{summarize-findings}x relative to k-epsilon, in the elliptic OpenFOAM solver on a wall-resolved mesh.}%
  \else\ifnum\pdfstrcmp{#1}{2}=0%
    \def\ckanswer@out{%
      Only partly, and we can say where it fails. The eigenframes of the stress anisotropy and the mean strain are misaligned by \ckvalue{misalignment\_pre}{44}{results/findings.json}{summarize-findings} degrees before transition and still by \ckvalue{misalignment\_turb}{24}{results/findings.json}{summarize-findings} degrees in equilibrium turbulence, so an eddy viscosity is structurally wrong before transition and imperfect everywhere. That bounds any eddy-viscosity closure, ours included.}%
  \else\ifnum\pdfstrcmp{#1}{3}=0%
    \def\ckanswer@out{%
      Load-bearing, so the hypothesis is refuted. Removing the rectifier raises the skin-friction error from \ckvalue{cf\_reference}{5.4\%}{results/findings.json}{summarize-findings} to \ckvalue{cf\_no\_clip}{73\%}{results/findings.json}{summarize-findings}, worse than k-epsilon. The lift-up term is load-bearing too, at \ckvalue{cf\_no\_liftup}{42\%}{results/findings.json}{summarize-findings} without it, while the streak viscous-decay term and the gate on omega production change the total score by \ckvalue{delta\_no\_decay}{-0.001}{results/findings.json}{summarize-findings} and \ckvalue{delta\_ungated}{0.000}{results/findings.json}{summarize-findings} and can be deleted outright.}%
  \else\ifnum\pdfstrcmp{#1}{4}=0%
    \def\ckanswer@out{%
      The vorticity Reynolds number Re\_v = y\textasciicircum{}2 Omega / nu. Every Re\_v variant beats every other driver tested, and the best structure is \ckvalue{best\_structure}{Rev\_p1.0\_a1}{results/findings.json}{summarize-findings}. The fitted threshold is \ckvalue{Lambda\_c}{418}{results/findings.json}{summarize-findings} against the classical value of 440, but the model is sensitive to it - using 440 itself raises the skin-friction error from \ckvalue{cf\_reference}{5.4\%}{results/findings.json}{summarize-findings} to \ckvalue{cf\_classical}{11.1\%}{results/findings.json}{summarize-findings} - and the fit wanders within search noise, which the paper quantifies, so the constant is not settled.}%
  \else\ifnum\pdfstrcmp{#1}{5}=0%
    \def\ckanswer@out{%
      Correlation, not redistribution. Through transition the u-v correlation coefficient grows \ckvalue{coherence\_ratio}{3.3}{results/findings.json}{summarize-findings}x while the energy partition changes by only \ckvalue{anisotropy\_ratio}{1.3}{results/findings.json}{summarize-findings}x; equivalently, the information shared between velocity components grows \ckvalue{total\_correlation\_ratio}{11}{results/findings.json}{summarize-findings}x.}%
  \else\ifnum\pdfstrcmp{#1}{6}=0%
    \def\ckanswer@out{%
      No. The component entropy is non-monotone through transition, falling to \ckvalue{entropy\_min}{0.50}{results/findings.json}{summarize-findings} at x = \ckvalue{entropy\_min\_x}{233}{results/findings.json}{summarize-findings} before recovering, so the constitutive relation is hysteretic rather than single-valued, and a closure built on it cannot hold the laminar state. Entropy is a good diagnostic and a poor substitute for the threshold.}%
  \else\ifnum\pdfstrcmp{#1}{7}=0%
    \def\ckanswer@out{%
      Yes, as a transient of transition rather than a standing inventory. Storage peaks at \ckvalue{peak\_storage}{24\%}{results/findings.json}{summarize-findings} of the energy routed into turbulence at x = \ckvalue{peak\_storage\_x}{246}{results/findings.json}{summarize-findings}, then collapses to a few percent. The turbulent share of entropy rejection rises from \ckvalue{rejection\_low}{0.520}{results/findings.json}{summarize-findings} to \ckvalue{rejection\_high}{0.559}{results/findings.json}{summarize-findings} with Reynolds number.}%
  \else\ifnum\pdfstrcmp{#1}{8}=0%
    \def\ckanswer@out{%
      Not yet. \ckvalue{n\_rectifier}{0}{results/findings.json}{summarize-findings} of the top eight evolved structures use the rectifier - an earlier answer that all eight did was falsified when the pipeline was re-run - and the stage does not reproduce from run to run, so the \ckvalue{n\_evaluated}{50}{results/findings.json}{summarize-findings} structures evaluated are evidence of nothing either way until the search is reproducible and properly powered.}%
  \else\ifnum\pdfstrcmp{#1}{9}=0%
    \def\ckanswer@out{%
      Yes. It reproduces the Blasius skin friction to within \ckvalue{cf\_error\_max}{1.3\%}{results/findings.json}{summarize-findings}, well below the differences between the models it ranks, and the bias is systematic so comparisons stay fair.}%
  \else\ifnum\pdfstrcmp{#1}{10}=0%
    \def\ckanswer@out{%
      Confirmed, and it splits the closure family. On the temporally evolving mixing layer of Lusher et al. (2026), the three mixing-length clip closures are indistinguishable from laminar, because their van Driest damping is keyed to a wall shear that does not exist, while \ckvalue{best\_closure}{clip-k-omega-gamma}{results/findings.json}{summarize-findings} is the best model on the case at \ckvalue{best\_score}{2.37}{results/findings.json}{summarize-findings}, ahead of Launder-Sharma at \ckvalue{ls\_score}{4.26}{results/findings.json}{summarize-findings}. Everything downstream of the plate needs a length scale made of transported quantities.}%
  \else\PackageError{calkit}{Unknown key '#1' for \string\ckanswer}{}%
  \fi\fi\fi\fi\fi\fi\fi\fi\fi\fi\fi%
  \ckanswer@out}%
\makeatother%
\makeatletter%
\newcommand\cknotes[1][all]{%
  \ifnum\pdfstrcmp{#1}{all}=0%
    \def\cknotes@out{%
      8, 11, 12, 13}%
  \else\ifnum\pdfstrcmp{#1}{8}=0%
    \def\cknotes@out{%
      The non-reproducibility is roadmap.md section 4, the highest-priority bug.}%
  \else\ifnum\pdfstrcmp{#1}{11}=0%
    \def\cknotes@out{%
      Open, and the most important gap. Everything here is calibrated on one DNS at one turbulence level, so the dependence cannot be seen at all. The fitted omega destruction coefficient of 0.1, against a standard 0.072, is a single-case compensation not to be trusted until a second flow tests it.}%
  \else\ifnum\pdfstrcmp{#1}{12}=0%
    \def\cknotes@out{%
      Open. Needs fluctuating velocity gradients, which means a new JHTDB pull. See docs/ideas-log.md section 3.}%
  \else\ifnum\pdfstrcmp{#1}{13}=0%
    \def\cknotes@out{%
      Open. They were written down but never tested in a forward solve. Two of them, a*grad(K) and b*(grad(U) dot grad(P)), are hard-coded into sim/newModel/solver/UEqn.H with fixed coefficients, which silently modifies the momentum equation for any model run through that solver.}%
  \else\PackageError{calkit}{Unknown key '#1' for \string\cknotes}{}%
  \fi\fi\fi\fi\fi%
  \cknotes@out}%
\makeatother%
\makeatletter%
\newcommand\ckevidence[1][all]{%
  \ifnum\pdfstrcmp{#1}{all}=0%
    \def\ckevidence@out{%
      1, 2, 3, 4, 5, 6, 7, 8, 9, 10}%
  \else\ifnum\pdfstrcmp{#1}{1}=0%
    \def\ckevidence@out{%
      \begin{itemize}\item value \texttt{results/findings.json}: \texttt{openfoam\_improvement\_over\_k\_epsilon}\item result \texttt{results/openfoam-vs-dns.json}\item figure \texttt{figures/openfoam-vs-dns.pdf}\item Section~\ref{sec:bench} (Benchmarking sensitivity)\end{itemize}}%
  \else\ifnum\pdfstrcmp{#1}{2}=0%
    \def\ckevidence@out{%
      \begin{itemize}\item value \texttt{results/findings.json}: \texttt{misalignment\_deg\_pretransition}\item value \texttt{results/findings.json}: \texttt{misalignment\_deg\_turbulent}\item result \texttt{results/flow-structure.json}\end{itemize}}%
  \else\ifnum\pdfstrcmp{#1}{3}=0%
    \def\ckevidence@out{%
      \begin{itemize}\item value \texttt{results/findings.json}: \texttt{ablation\_reference\_cf\_rel\_rms}\item value \texttt{results/findings.json}: \texttt{ablation\_no\_clip\_cf\_rel\_rms}\item value \texttt{results/findings.json}: \texttt{ablation\_no\_liftup\_cf\_rel\_rms}\item value \texttt{results/findings.json}: \texttt{ablation\_no\_streak\_decay\_delta}\item value \texttt{results/findings.json}: \texttt{ablation\_ungated\_omega\_delta}\item result \texttt{results/closure-ablation.json}\end{itemize}}%
  \else\ifnum\pdfstrcmp{#1}{4}=0%
    \def\ckevidence@out{%
      \begin{itemize}\item value \texttt{results/findings.json}: \texttt{best\_structure}\item value \texttt{results/findings.json}: \texttt{fitted\_threshold\_Lambda\_c}\item value \texttt{results/findings.json}: \texttt{ablation\_reference\_cf\_rel\_rms}\item value \texttt{results/findings.json}: \texttt{ablation\_classical\_threshold\_cf\_rel\_rms}\item result \texttt{results/closure-search.json}\item Section~\ref{sec:fails} (What does not transfer)\end{itemize}}%
  \else\ifnum\pdfstrcmp{#1}{5}=0%
    \def\ckevidence@out{%
      \begin{itemize}\item value \texttt{results/findings.json}: \texttt{coherence\_ratio}\item value \texttt{results/findings.json}: \texttt{anisotropy\_ratio}\item value \texttt{results/findings.json}: \texttt{total\_correlation\_ratio}\item result \texttt{results/flow-structure.json}\end{itemize}}%
  \else\ifnum\pdfstrcmp{#1}{6}=0%
    \def\ckevidence@out{%
      \begin{itemize}\item value \texttt{results/findings.json}: \texttt{entropy\_min}\item value \texttt{results/findings.json}: \texttt{entropy\_min\_x}\item result \texttt{results/flow-structure.json}\end{itemize}}%
  \else\ifnum\pdfstrcmp{#1}{7}=0%
    \def\ckevidence@out{%
      \begin{itemize}\item value \texttt{results/findings.json}: \texttt{peak\_turbulent\_storage\_fraction}\item value \texttt{results/findings.json}: \texttt{peak\_turbulent\_storage\_x}\item value \texttt{results/findings.json}: \texttt{turbulent\_rejection\_fraction\_low\_Re}\item value \texttt{results/findings.json}: \texttt{turbulent\_rejection\_fraction\_high\_Re}\item result \texttt{results/exergy-budget.json}\item Section~\ref{sec:transfers} (What transfers)\end{itemize}}%
  \else\ifnum\pdfstrcmp{#1}{8}=0%
    \def\ckevidence@out{%
      \begin{itemize}\item value \texttt{results/findings.json}: \texttt{evolved\_top8\_using\_rectifier}\item value \texttt{results/findings.json}: \texttt{evolved\_structures\_evaluated}\item result \texttt{results/closure-evolution.json}\end{itemize}}%
  \else\ifnum\pdfstrcmp{#1}{9}=0%
    \def\ckevidence@out{%
      \begin{itemize}\item value \texttt{results/findings.json}: \texttt{screening\_solver\_cf\_error\_max}\item result \texttt{results/blasius-validation.json}\item Section~\ref{sec:method} (Method)\end{itemize}}%
  \else\ifnum\pdfstrcmp{#1}{10}=0%
    \def\ckevidence@out{%
      \begin{itemize}\item value \texttt{results/findings.json}: \texttt{mixing\_layer\_closures\_identical\_to\_laminar}\item value \texttt{results/findings.json}: \texttt{mixing\_layer\_best\_closure}\item value \texttt{results/findings.json}: \texttt{mixing\_layer\_best\_normalized}\item value \texttt{results/findings.json}: \texttt{mixing\_layer\_launder\_sharma\_normalized}\item result \texttt{results/benchmark.json}\item Section~\ref{sec:scoring} (Scoring across cases)\end{itemize}}%
  \else\PackageError{calkit}{Unknown key '#1' for \string\ckevidence}{}%
  \fi\fi\fi\fi\fi\fi\fi\fi\fi\fi\fi%
  \ckevidence@out}%
\makeatother%
\newcommand\ckfindings{\par\noindent\textbf{Q1. \ckquestion[1]}\ckblock{1}{calkit.yaml}\par\noindent\ckanswer[1]\ckevidence[1]\par
\par\noindent\textbf{Q2. \ckquestion[2]}\ckblock{2}{calkit.yaml}\par\noindent\ckanswer[2]\ckevidence[2]\par
\par\noindent\textbf{Q3. \ckquestion[3]}\ckblock{3}{calkit.yaml}\par\noindent\ckanswer[3]\ckevidence[3]\par
\par\noindent\textbf{Q4. \ckquestion[4]}\ckblock{4}{calkit.yaml}\par\noindent\ckanswer[4]\ckevidence[4]\par
\par\noindent\textbf{Q5. \ckquestion[5]}\ckblock{5}{calkit.yaml}\par\noindent\ckanswer[5]\ckevidence[5]\par
\par\noindent\textbf{Q6. \ckquestion[6]}\ckblock{6}{calkit.yaml}\par\noindent\ckanswer[6]\ckevidence[6]\par
\par\noindent\textbf{Q7. \ckquestion[7]}\ckblock{7}{calkit.yaml}\par\noindent\ckanswer[7]\ckevidence[7]\par
\par\noindent\textbf{Q8. \ckquestion[8]}\ckblock{8}{calkit.yaml}\par\noindent\ckanswer[8]\ckevidence[8]\par
\par\noindent\textbf{Q9. \ckquestion[9]}\ckblock{9}{calkit.yaml}\par\noindent\ckanswer[9]\ckevidence[9]\par
\par\noindent\textbf{Q10. \ckquestion[10]}\ckblock{10}{calkit.yaml}\par\noindent\ckanswer[10]\ckevidence[10]\par}%


\begin{document}
\maketitle

\begin{abstract}
Every number in this draft is produced by the pipeline in the accompanying
repository. DRAFT --- written by Claude (Anthropic's Claude Code) at the
author's request as a fact-carrying skeleton, to be rewritten.
\end{abstract}

\section{Introduction}\label{sec:introduction}

Data-driven turbulence closure development \citep{Duraisamy2019} is usually
reported through in-sample agreement with the data used to build it. We
develop a transition closure for a bypass-transitional flat-plate boundary
layer, and find that the standard diagnostics do not distinguish a
constitutive law from a curve fit. Three independent routes --- a-posteriori
structure search, a-priori term regression, and an integral flux budget ---
each produce excellent in-sample agreement and each fails to transfer.

We report the failure modes quantitatively, because each is invisible in the
metrics normally quoted, and each inflated our own results before it was
found.

\section{Case and data}\label{sec:case}

The primary dataset is the JHTDB \texttt{transition\_bl} DNS
\citep{JHTDBDescription}: a flat plate under free-stream turbulence,
$x \in [\result[XInlet], \result[XOutlet]]$, $y \in [0, \result[YMax]]$, $\nu = 1.25\times10^{-3}$, with
$Tu = \result[TuInlet]\,\%$ and $\delta_{99} = \result[DeltaInlet]$ already developed at the inlet
station. Skin friction falls to a minimum at $x \approx 205$ and peaks near
$x \approx 450$. An independent zero-pressure-gradient turbulent boundary
layer DNS at $Re_\theta = 4000$--$6500$ \citep{Sillero2013} is used
throughout as an out-of-sample test.

\section{Method}\label{sec:method}

\subsection{Solvers}

Two fidelities are used. Structure and coefficient searches are screened with
a parabolic boundary-layer marching solver, which advances the mean momentum
and closure equations downstream from the DNS inlet profile in about a second
per candidate. Every result quoted for a model is then reproduced in
OpenFOAM v2412 (\texttt{simpleFoam}, steady SIMPLE) on a wall-resolved mesh,
because the parabolic solver is optimistic: it is handed the inlet profile,
marches without an elliptic pressure field, and returned roughly a factor of
three lower skin-friction error than the elliptic solve for the same closure
and coefficients.

The elliptic domain is a single block matched to the DNS extent,
$x \in [\result[XInlet], \result[XOutlet]]$ and $y \in [0, \result[YMax]]$,
with 80 cells in $y$ graded to a first cell at $y^+ \lesssim 1$. Placing the
inlet at the DNS station rather than upstream of a leading edge is not
cosmetic; see \S\ref{sec:bench}.

\subsection{Inlet conditions}

Every model is given the measured DNS profile of $U$ and $k$ at the inlet
station. The inlet $\omega$ cannot be shared, because under the analytic
free-decay law $k \propto (1 + \beta \omega_0 (x-x_{\rm in})/U)^{-\beta^*/\beta}$
the decay exponent is $\beta^*/\beta$; a single $\omega_0$ would give models
with different $\beta$ different free-stream decay. It is therefore fitted per
model to the same measured decay just downstream of the inlet. For the
$\gamma$--$Re_{\theta t}$ model the inlet $Re_{\theta t}$ is taken from that
model's own empirical correlation at the inlet turbulence intensity, giving
$\result[ReThetatInlet]$.

\subsection{Coefficient fitting}

The closure under test transports $k$, $\omega$ and an activation fraction
$\gamma$, with the eddy viscosity gated as $\nu_t = \gamma k/\omega$ and the
$\gamma$ source rectified about a threshold in the vorticity Reynolds number.
It has \result[NFittedCoeffs]{} fitted coefficients.

Fitting is by random search over the coefficient bounds
(\result[NRandom]{} samples) followed by \result[NRefineRounds]{} rounds of
Gaussian refinement about the incumbent (\result[NRefine]{} samples each, with
successively shrinking width). The objective is a weighted sum of relative
errors in skin friction, mean velocity, momentum thickness, free-stream decay
and turbulence energy, each divided by a target value so that a contribution
of one means that quantity has arrived. The turbulence-energy term is included
because an objective built only on the mean field cannot see a closure that
carries an order of magnitude too little $k$.

\subsection{A-priori regression}

Candidate closure terms are evaluated directly from the DNS and their
coefficients solved for by least squares, rather than being scored after a
forward solve. The target is normalised by $k$, making it the structure
parameter and the library dimensionless; an unnormalised target spans orders
of magnitude across the plate and the fit would see only the turbulent end.

Two guards are applied. The library admits only quantities a RANS solver
possesses, so that no column carries information from the target, and the
correlation of every column with the target is reported so that a
near-identity is visible rather than being read as a result. Fits are then
evaluated on data not used to produce them: the opposite half of the plate,
and the independent DNS.

\subsection{Reproducibility}

Every number and figure in this paper is produced by a pipeline stage from
version-controlled inputs, and the values quoted in the text are generated
from the results files rather than transcribed.

\section{What transfers}\label{sec:transfers}

Two measured quantities survive cross-validation between the two DNS.

The structure parameter $a_1 = -\langle u'v'\rangle / 2k$, integrated across
the layer, is $\result[aOneTransitional]$ at the downstream end of the transitional case and
$\result[aOneJimenez] \pm \result[aOneJimenezSd]$ across $Re_\theta = 4000$--$6500$ in the independent DNS
--- agreement to \result[aOneRelDiff]\,\% across a fourfold change in $Re_\theta$. This is
Bradshaw's structure parameter \citep{Bradshaw1967} and is not new; it is
reported here because it is the only quantity in this study that behaves like
a constant.

The dissipation coefficient $C_\epsilon = \epsilon L / u'^3$, with
$L = \delta_{99}$ and $u' = \sqrt{2\bar{k}/3}$, is obtained without modelling
from the exact integral balance
\begin{equation}
  \epsilon_{\rm int} = P_{\rm int} - \frac{{\rm d}}{{\rm d}x}
  \int U k \, {\rm d}y ,
  \label{eq:intbudget}
\end{equation}
in which the transport terms are internal and integrate to vanishing boundary
fluxes. It is $\result[CepsTurb] \pm \result[CepsTurbSd]$ over $Re_\theta = 900$--$1500$ in the
transitional case and $\result[CepsJimenez]$ in the independent DNS.
 Through transition it varies by a factor of \result[CepsRatio],
rising from $\result[CepsPre]$; the ratio $P/\epsilon$ peaks at $\result[PoverEpsPeak]$ near
$x \approx \result[PoverEpsPeakX]$, coinciding with the peak in turbulent energy storage
($x \approx 246$) and the minimum of the component entropy of the fluctuation
field ($x \approx 233$). That the cascade is out of equilibrium through
transition is established \citep{Vassilicos2015, LiuFangFang2021}; the
measurement is included as validation of the budget, not as a new result.

\begin{figure}
  \centerline{\ckfigure[width=0.9\textwidth]{../figures/dissipation.pdf}}
  \caption{Left: dissipation coefficient, with the equilibrium band of this
    DNS shaded and the independent DNS dashed. Right: the flux imbalance,
    peaking at transition onset.}
  \label{fig:dissipation}
\end{figure}


\section{What does not transfer}\label{sec:fails}

\subsection{A-priori term regression}

Regressing $-\langle u'v'\rangle/k$ on a library of \result[NTerms] dimensionless groups
formed only from quantities a RANS solver possesses ($\partial U/\partial y$,
$\partial^2 U/\partial y^2$, $y$, $\nu$, $k$) gives $R^2 = \result[StressInSample]$ in sample.
Coefficients fitted downstream of $x = 450$ and applied upstream give
$R^2 = \result[StressDownToUp]$; coefficients fitted on the transitional DNS and applied to the
independent DNS give $R^2 = \result[StressToJimenez]$, although the latter fits itself to
$R^2 = \result[JimenezSelfFit]$. The normalised disagreement between upstream and downstream
coefficient vectors is $\result[StressCoeffDisagree]$, its maximum. The same holds for the unclosed
residual of the $k$ budget: $R^2 = \result[EpsInSample]$ in sample, $\result[EpsUpToDown]$ and $\result[EpsDownToUp]$ out.

Every regime is separately and easily fittable; no coefficient set serves two
of them.
Table~\ref{tab:transfer} collects the splits and figure~\ref{fig:transfer}
shows the in-sample against out-of-sample contrast.

\begin{table}
  \begin{center}
  \begin{tabular}{lcc}
    split & $-\langle u'v'\rangle/k$ & $k$-budget residual \\[3pt]
    in sample                         & \result[StressInSample]  & \result[EpsInSample] \\
    fit upstream, predict downstream  & ---                      & \result[EpsUpToDown] \\
    fit downstream, predict upstream  & \result[StressDownToUp]  & \result[EpsDownToUp] \\
    fit here, predict independent DNS & \result[StressToJimenez] & --- \\
    independent DNS fitted on itself  & \result[JimenezSelfFit]  & --- \\
  \end{tabular}
  \caption{Coefficient of determination for the a-priori regression with
    \result[NTerms]{} admissible terms. Negative values are worse than
    predicting the mean.}
  \label{tab:transfer}
  \end{center}
\end{table}

\begin{figure}
  \centerline{\ckfigure[width=0.5\textwidth]{../figures/transfer.pdf}}
  \caption{The same library fits either target in sample and fails out of it.}
  \label{fig:transfer}
\end{figure}


\subsection{Trivial solutions}

Two failure modes produce high scores with no content.

A library column carrying information from the target reconstructs it. With
$-\langle u'v'\rangle = R_{uv} u' v'$ and $u'^2/2k$ approximately steady,
$-\langle u'v'\rangle/k$ is nearly proportional to $v'/u'$, so including
$v'/u'$ --- unavailable to a RANS solver in any case --- raised $R^2$ from
$0.361$ to $0.965$ on an impoverished four-term library. On the \result[NTerms]-term
library the same addition moves $\result[StressInSample]$ to $\result[StressDiagnostic]$, showing the first result
measured library poverty rather than missing physics.

Monotonic variables correlate with monotonic targets. $C_\epsilon$ correlates
with the activation $\gamma = a_1/a_{1\infty}$ at $r = \result[CollapseGammaCorr]$, yet collapses
onto it no better than $\result[CollapseGamma]\,\%$ relative RMS --- worse than onto the
streamwise coordinate $x$ itself ($\result[CollapseX]\,\%$), which carries no physical
content. Correlation coefficients are not evidence of a closure relation; a
trivial-coordinate baseline should be reported alongside any claimed collapse.
\begin{figure}
  \centerline{\ckfigure[width=0.55\textwidth]{../figures/collapse.pdf}}
  \caption{Collapse error for each candidate against the streamwise coordinate
    baseline (dashed), with correlation coefficients annotated. The
    best-correlated variable is not the best-collapsing one.}
  \label{fig:collapse}
\end{figure}


\subsection{Structure selection inside search noise}

Model structures are commonly ranked by the best error their inner
coefficient fit reaches. Repeating that inner fit with an identical protocol
and only the random seed changed gives a standard deviation of $\result[FitNoiseSd]$ in the
objective (\result[NSeeds] seeds, \result[NVariants] structural variants), against between-structure
gaps of $0.2$--$0.6$. The fitted transition threshold $\Lambda_c$ wanders
between $\result[LamCMin]$ and $\result[LamCMax]$ for one fixed structure. No structural conclusion in
this study is supported by its score, and the same objection applies to any
outer-loop structure search whose inner fit is not first shown to be
converged.
Table~\ref{tab:noise} and figure~\ref{fig:noise} give the spread.

\begin{table}
  \begin{center}
  \begin{tabular}{lcccc}
    variant & mean & s.d. & min & max \\[3pt]
    komega+gate & \result[NoiseKomegaGateMean] & \result[NoiseKomegaGateSd] & \result[NoiseKomegaGateMin] & \result[NoiseKomegaGateMax] \\
    mixing      & \result[NoiseMixingMean]     & \result[NoiseMixingSd]     & \result[NoiseMixingMin]     & \result[NoiseMixingMax] \\
    mixing+gate & \result[NoiseMixingGateMean] & \result[NoiseMixingGateSd] & \result[NoiseMixingGateMin] & \result[NoiseMixingGateMax] \\
  \end{tabular}
  \caption{Objective across \result[NSeeds]{} random seeds per structural
    variant under an identical protocol. Between-variant differences in the
    mean are smaller than the within-variant spread.}
  \label{tab:noise}
  \end{center}
\end{table}

\begin{figure}
  \centerline{\ckfigure[width=0.55\textwidth]{../figures/fit-noise.pdf}}
  \caption{Each marker is one seed; bars span one standard deviation.}
  \label{fig:noise}
\end{figure}


\section{Benchmarking sensitivity}\label{sec:bench}

Transition-model comparisons are dominated by case setup. The
$\gamma$--$Re_{\theta t}$ model \citep{LangtryMenter2009} was run on the same
wall-resolved mesh under two configurations differing in no model coefficient:

\begin{table}
  \begin{center}
  \begin{tabular}{lcc}
    configuration & $U$ rel. RMS & $c_f$ rel. err. \\[3pt]
    inlet 127 units upstream, $Re_{\theta t}$ set arbitrarily & \result[LMOldU] & \result[LMOldCf] \\
    inlet at the DNS station, $Re_{\theta t} = \result[ReThetatInlet]$ from its own correlation & \result[LMNewU] & \result[LMNewCf] \\
  \end{tabular}
  \caption{Effect of case setup alone on a transition model.}
  \label{tab:setup}
  \end{center}
\end{table}

The upstream run-up is the larger effect. The DNS begins at the plate with a
developed boundary layer; a domain extending $127$ units further upstream
requires the free-stream turbulence to be older at the DNS inlet station than
the data allows, and matching the measured level then demands an inlet
intensity of $20.5\,\%$, which trips the model at the leading edge. Placing
the inlet at the DNS station gives $Tu = \result[TuInlet]\,\%$ for every model.

On the corrected domain, with inlet conditions fitted per model to the same
measured free-stream decay, $\gamma$--$Re_{\theta t}$ attains $c_f$ error
$\result[LMNewCf]$ using published coefficients, against $\result[ClipNewCf]$ for the present
closure using eight coefficients calibrated on this same case.

\begin{table}
  \begin{center}
  \begin{tabular}{lccc}
    model & $U$ rel. RMS & $c_f$ mean & $c_f$ max \\[3pt]
    $\gamma$--$Re_{\theta t}$ & \result[DnsKOmegaSstLmU] & \result[DnsKOmegaSstLmCf] & \result[DnsKOmegaSstLmCfMax] \\
    present closure           & \result[DnsClipKGammaU]  & \result[DnsClipKGammaCf]  & \result[DnsClipKGammaCfMax] \\
    $k$--$\omega$ SST         & \result[DnsKOmegaSstU]   & \result[DnsKOmegaSstCf]   & \result[DnsKOmegaSstCfMax] \\
    $k_L$--$k_T$--$\omega$    & \result[DnsKklOmegaU]    & \result[DnsKklOmegaCf]    & \result[DnsKklOmegaCfMax] \\
    laminar                   & \result[DnsLaminarU]     & \result[DnsLaminarCf]     & \result[DnsLaminarCfMax] \\
  \end{tabular}
  \caption{All models on the domain matched to the DNS, same mesh, with inlet
    conditions fitted per model to the same measured free-stream decay.}
  \label{tab:models}
  \end{center}
\end{table}

\begin{figure}
  \centerline{\ckfigure[width=0.7\textwidth]{../figures/model-comparison.pdf}}
  \caption{Skin friction against the DNS (markers) on the matched domain.}
  \label{fig:benchmark}
\end{figure}

\subsection{Scoring across cases}\label{sec:scoring}

% TODO (human): write this subsection. Skeleton only; the definitions below
% are what scripts/run-benchmark.py and pypkg/cases/base.py implement, so the
% text must match them. Points the text has to make (docs/rans-gym.md
% sections 1 and 3 give the reasoning behind each):
%   - every case declares a target error per metric, the error at which a
%     careful reader would call the model and the data indistinguishable;
%   - the score is the mean over metrics of error/target, so 1.0 means the
%     same thing on a transitional plate, a channel and a mixing layer;
%   - every closure declares the cases its coefficients were fitted on, and
%     in-sample and out-of-sample scores are reported separately and never
%     averaged together (the whole point, given Sec. \ref{sec:fails});
%   - every closure is seeded by the same case-supplied rule on cases whose
%     inlet is already turbulent, and the seeded station or transient is
%     excluded from scoring;
%   - a diverged or non-converged run scores infinity, not a number;
%   - a peak the model can make grid-dependent is reported but not scored.
% Cite Sansica2026 for the initial-condition point (their Sec. 3.1).

A case $c$ declares, for each metric $m$ it scores, a target error
$\tau_{c,m}$; the normalized score of closure $\mathcal{M}$ on case $c$ is
%
\begin{equation}
  s_c(\mathcal{M}) = \frac{1}{|M_c|} \sum_{m \in M_c}
    \frac{e_{c,m}(\mathcal{M})}{\tau_{c,m}},
  \label{eq:score}
\end{equation}
%
so $s_c = 1$ means the closure matches the data to within the declared target
on every metric. Each closure declares the set $\mathcal{C}(\mathcal{M})$ of
cases its coefficients were fitted on, and the leaderboard reports
%
\begin{equation}
  \bar{s}_{\mathrm{in}} = \langle s_c \rangle_{c \in \mathcal{C}(\mathcal{M})},
  \qquad
  \bar{s}_{\mathrm{out}} = \langle s_c \rangle_{c \notin \mathcal{C}(\mathcal{M})},
  \qquad
  \pi = \bar{s}_{\mathrm{out}} / \bar{s}_{\mathrm{in}},
  \label{eq:split}
\end{equation}
%
with the transfer penalty $\pi$ the number this paper is about. A run that
diverges or fails to converge scores $s_c = \infty$.


\section{Conclusions}\label{sec:conc}

In-sample agreement, correlation coefficients, and structure rankings all
failed to distinguish a transferable relation from a case-specific fit here.
The diagnostics that did discriminate were out-of-sample transfer to an
independent DNS, a trivial-coordinate baseline for claimed collapses, and
repetition of the inner fit under different random seeds. We recommend that
data-driven closure studies report all three.


\bibliographystyle{jfm}
\bibliography{../references}

\appendix
\section{Questions and answers}\label{sec:questions}

% Generated from the questions in calkit.yaml by the questions-to-latex
% stage: the project's record of what it asked and found, rendered here from
% the same results files the text above cites.
\ckfindings

\end{document}
